\begin{table}[t]\centering\small
\begin{tabular}{lrr}
\toprule
sectors with a nonzero diagonal & random rings & share \\
\midrule
$\{\Oz,\OEc\}$ & 126 & 63.0\,\% \\
$\{\Oz,\OEc,\ODtwo\}$ & 47 & 23.5\,\% \\
$\{\OEc\}$ & 19 & 9.5\,\% \\
$\{\Oz,\OEc,\ODthree\}$ & 2 & 1.0\,\% \\
$\{\Oz,\OEc,\ODthree,\ODtwo\}$ & 2 & 1.0\,\% \\
$\{\OEc,\ODtwo\}$ & 2 & 1.0\,\% \\
$\{\Oz\}$ & 1 & 0.5\,\% \\
$\{\Oz,\OEc,\OTtwo\}$ & 1 & 0.5\,\% \\
\midrule
$J_5$, $J_7$, $J_9$ (rebuilt from claw words) & \multicolumn{2}{r}{$\{\OTone\}$ in each} \\
order of the random rings & \multicolumn{2}{r}{mean $62.6 \pm 25.4$, range 18 to 132} \\
\bottomrule
\end{tabular}
\caption{Sector traces of the 200 random rings of the cross-check (computation EXP-014, \texttt{crosscheck.json}): for each set of sectors, the number of rings whose product of sector matrices has a nonzero diagonal exactly in those sectors. In every case the solver also found a Petersen coloring. The $\pm$ value is the sample standard deviation of the orders.}\label{tab:cross}
\end{table}
